\color{unidarkgreen}{\subsection{Токовая защита нулевой последовательности резистора заземления нейтрали}}
\color{black}

\subsubsection{Общие сведения}

\begin{enumerate}[label=\arabic{section}.\arabic{subsection}.\arabic{subsubsection}.\arabic*,  align=left, labelsep=4pt, leftmargin=0pt, itemindent=57pt]

\item
Функциональный блок ТЗНП РЗН предназначен для защиты РЗН от замыканий на землю в сетях
с резистивным заземлением нейтрали. ТЗНП РЗН подключается к трансформатору тока на вводе
резистора заземления нейтрали силовых трансформаторов, работающих в режиме низкоомного
заземления нейтрали стороны. Защита выполняется двухступенчатой и действует как резервная с
выдержкой времени на отключение ввода НН и на отключение трансформатора.

\item
Ступень состоит из максимального органа сравнения с уставкой \textbf{«Iср ТЗНП РЗН1»} и выдержки
времени \textbf{««Тср ТЗНП РЗН1»}. Предусмотрены внешние сигналы вывода ступени из работы и блокировки,
а также программный переключатель ввода/вывода ступени функционального блока «ТЗНП РЗН1»
(рисунок 17). Ступень ТЗНП РЗН может быть заблокирована от блока детектора второй гармоники,
который вводится программным переключателем \textbf{««БНТ ТЗНП РЗН1»}, при превышении уровня второй
гармоники тока Iрзн уставки \textbf{««I2г/I1г ТЗНП РЗН1»} (рисунок \ref{pic:TZNP_RZN}). Ввод/вывод блокировки по второй
гармонике для каждой ступени производится программным переключателем \textbf{««Реж БНТ ТЗНП РЗН1»}.

\item
Логика «ТЗНП РЗН2» аналогична «ТЗНП РЗН1».

\begin{figure}[H]
\centering
\includegraphics[width=0.7\textwidth,height=0.7\textheight,keepaspectratio]{\fbpath/031.0151 ТЗНП РЗН/img/TZNP_RZN.pdf}
\captionof{figure}{Функциональный блок ТЗНП РЗН}\label{pic:TZNP_RZN}
\end{figure}


\small
\begin{longtable}{|>{\centering\arraybackslash}m{5.3cm}|>{\centering\arraybackslash}m{4.2cm}|>{\centering\arraybackslash}m{1cm}|>{\centering\arraybackslash}m{3.3cm}|>{\centering\arraybackslash}m{1.8cm}|}
% Первая страница: полный заголовок
\caption{Перечень уставочных параметров ТЗНП ВН\hfill\vspace{-0.5\baselineskip}}\label{desc:tbl1}\\
\hline
\rowcolor{gray!30}
Наименование уставки & Диапазон & Шаг & Значение по умолчанию & Единица измерения \\
\hline
\endfirsthead

% Промежуточные страницы: «Продолжение таблицы»
\caption*{\hspace{3pt}\emph{Продолжение таблицы \ref{desc:tbl1}}\hfill\ }\vspace{-0.5\baselineskip} \\
\hline
\rowcolor{gray!30}
Наименование уставки & Диапазон & Шаг & Значение по умолчанию & Единица измерения \\
\hline
\endhead

%===m>ДТЗ|-
%===m
\end{longtable}
\normalsize	

\end{enumerate}